% Created 2026-04-11 Sat 18:59
% Intended LaTeX compiler: pdflatex
\documentclass[11pt]{article}
\usepackage[utf8]{inputenc}
\usepackage[T1]{fontenc}
\usepackage{graphicx}
\usepackage{longtable}
\usepackage{wrapfig}
\usepackage{rotating}
\usepackage[normalem]{ulem}
\usepackage{amsmath}
\usepackage{amssymb}
\usepackage{capt-of}
\usepackage{hyperref}
\usepackage[left=0.75in,right=0.75in,top=1in,bottom=1in]{geometry}
\usepackage{setspace}\n\doublespacing
\usepackage{enumitem}
\setlist[itemize]{itemsep=2pt, topsep=4pt}
\author{Ben Heinze}
\date{\today}
\title{Final Project - Part 2}
\hypersetup{
 pdfauthor={Ben Heinze},
 pdftitle={Final Project - Part 2},
 pdfkeywords={},
 pdfsubject={},
 pdfcreator={Emacs 29.3 (Org mode 9.6.15)}, 
 pdflang={English}}
\begin{document}

\maketitle
\tableofcontents



\section{Instructions}
\label{sec:orgf73a3ce}

Read Higgs and Amhrien and revisit Ch 23/24 as necessary. Turn in a document that addresses the following.

\begin{enumerate}
\item Read Higgs and Amhrein and Chp 23.4 and 23.5. In your own words, translate the cautions presented in Amhrein and Higgs to the content presented in 23.4 and 23.5. That is what "baggage" comes with the powre calculation methods discussed in the Sleth (Max 1/2 page double spaced).

\item \textbf{\textbf{Quantitative Backdrop}} Draw a "Quantitative Backdrop" for the parameter associated with a component of your reserach question that you could summarize with a confidence interval. Provide a complete picture with labeled intervals and complete descriptions of the practical implications of each interval.

\item Choose \emph{one} of the intervals in your backdrop (e.g., an interval that spans gray area and practically meaningful results). Write a \emph{hypothetical} summary of statistical findings to communicate results from this interval. Be sure to contextualize with appropriate scope of inference and practical context from your quantitative background.
\end{enumerate}
\end{document}
